\begin{description}
\item[(T10.1)] The ratio of the number density of protons to that of
  neutrons at $t_{\mathrm{wk}}$ will be given by
  \[ \frac{N_{\mathrm{n}}(t_{\mathrm{wk}})}{N_{\mathrm{p}}(t_{\mathrm{wk}})}
     = \left(\frac{m_{\mathrm{n}}}{m_{\mathrm{p}}}\right)^{3/2}
       \exp\left(-\frac{(m_{\mathrm{n}} - m_{\mathrm{p}})c^2}
                        {k_{\mathrm{B}}T_{\mathrm{wk}}}\right) \]
  \hfill(2.0)

  Substituting the values of the masses of neutron and proton, and
  $k_{\mathrm{B}}T_{\mathrm{wk}} = 800$\,keV, one obtains
  \[ X_{\mathrm{n}}(t_{\mathrm{wk}})
     = \frac{N_{\mathrm{n}}}{N_{\mathrm{p}} + N_{\mathrm{n}}} = 0.166 \]
  \hfill(2.0)

\item[(T10.2a)] For a black body spectrum,
  $\lambda_{\mathrm{max}}T$ should be equal to a constant. \hfill(1.0)

  Since the wavelength of light is $\lambda \propto a$, this implies
  $T \propto a^{-1}$, which implies $\beta = -1$. \hfill(1.0)

  Alternatively, for black body radiation, energy density is proportional
  to $T^4$, but the energy density is proportional to $a^{-4}$.
  Therefore, $T \propto a^{-1}$.

\item[(T10.2b)] Graph B should be ticked. Wien's law removes options D
  and E, Planck curves at different temperatures should not intersect,
  and so C and A are also incorrect.

\item[(T10.3a)] Each of the 4 points correctly located on the graph
  \hfill(1.0)

  Drawing a smooth curve passing through all the points (within one grid
  point) \hfill(1.0)

  Non-smooth lines, or straight lines not passing through all points get
  0 credit.
  % FIGURE NEEDED: solution plot of N_D/N_n (log scale 10^-2..10^2)
  % versus k_B T = 58..72 keV with 4 marked points and smooth declining
  % curve (2025_theory_solution.pdf page 38)

\item[(T10.3b)] The value of
  $k_{\mathrm{B}}T_{\mathrm{nuc}} = 66.0\,\mathrm{keV}$. \hfill(1.0)

  Full credit range: 65.4\,keV to 66.6\,keV.

\item[(T10.3c)] As $T \propto a^{-1}$, and $a \propto t^{1/2}$, we have
  $T \propto t^{-1/2}$. \hfill(2.0)

  The temperature and time when all neutrons and protons fuse into
  Helium are $T_{\mathrm{nuc}}$ (66.0\,keV) and $t_{\mathrm{nuc}}$,
  respectively.
  \[ t_{\mathrm{nuc}} = t_{\mathrm{wk}}
     \left(\frac{T_{\mathrm{wk}}}{T_{\mathrm{nuc}}}\right)^2
     = 250\,\mathrm{s} \]
  \hfill(2.0)

\item[(T10.4)] During $t_{\mathrm{wk}} < t < t_{\mathrm{nuc}}$, as the
  Universe cools from 800\,keV to 66.0\,keV, neutrons decay into protons
  and with a half life of 610.4\,s, which corresponds to a mean life time
  $\tau = 880.6$\,s. \hfill(1.0)

  The quantity $(N_{\mathrm{p}} + N_{\mathrm{n}})a^3$ does not change as
  the increase in protons is compensated by the decrease in neutrons.
  \hfill(1.0)

  However, $N_{\mathrm{n}}a^3$ will go down by an exponential factor
  compared to its value computed in the first question as the neutrons
  decay into protons, i.e.,
  \[ X_{\mathrm{n,\,nuc}} = X_{\mathrm{n,\,wk}}
     \exp(-(t_{\mathrm{nuc}} - t_{\mathrm{wk}}/880.6)) \]
  % sic: misplaced parenthesis in source; the numerical line below has
  % the intended (t_nuc - t_wk)/880.6
  \hfill(2.0)
  \[ X_{\mathrm{n,\,nuc}} = X_{\mathrm{n,\,wk}}
     \exp(-(250 - 1.70)/880.6) = 0.125 \]
  \hfill(1.0)

  Give part credit of 1 marks if half life time is used instead of the
  mean life time in the denominator of the exponential.

\item[(T10.5)] Helium forms with the combination of 2 protons and 2
  neutrons (ignoring mass deficit), and thus a number density
  $N_{\mathrm{n}}$ of neutrons and correspondingly same number density of
  protons will be locked in Helium. The number density of both species
  goes as $a^{-3}$, and so will cancel from the numerator and
  denominator. \hfill(1.0)

  Even after Helium is formed the sum of the mass of Helium and Hydrogen
  will still be the mass of all the neutrons and protons. \hfill(1.0)

  Thus, the abundance of Helium by mass should be
  \[ Y_{\mathrm{prim}}
     = \frac{(m_{\mathrm{p}} + m_{\mathrm{n}})N_{\mathrm{n}}}
            {m_{\mathrm{p}}N_{\mathrm{p}} + m_{\mathrm{n}}N_{\mathrm{n}}}
     \approx 2X_{\mathrm{n,\,nuc}} = 0.250 \]
  \hfill(1.0)

  Alternatively, the value of $X_{\mathrm{n,\,nuc}}$ is $1/8$. After
  formation of Helium all neutrons will combine with protons. Therefore,
  the mass in Helium will be $2/8$ of the total mass.

\item[(T10.6a)] By fitting a line to these data points the Helium
  abundance for the said galaxy would be: $Y = 0.2480$ \hfill(2.0)

  Full credit range: 0.2475 to 0.2485.
  % FIGURE NEEDED: annotated straight-line fit through the Y vs (O/H)
  % data, marked points (20, 0.2445) and (195, 0.2485), slope = 22.86
  % (2025_theory_solution.pdf page 40)

\item[(T10.6b)] The slope of the line yields
  $dY/d(\mathrm{O/H}) = 23$. \hfill(2.0)

  Full credit range: 19 to 27.

\item[(T10.6c)] The primordial Helium abundance,
  $Y_{\mathrm{prim}}^{\mathrm{obs}} = 0.2440$. \hfill(2.0)

  Full credit range: 0.2435 to 0.2445.

\item[(T10.7)] From the Saha equation, we see that
  $N_{\mathrm{D}}/N_{\mathrm{n}}(T) \propto \eta$, thus when
  $\eta\downarrow$, $N_{\mathrm{D}}/N_{\mathrm{n}}(T)$ will also
  $\downarrow$. From the plot in (T10.3a), this would mean that
  $N_{\mathrm{D}}/N_{\mathrm{n}}(T)$ would be unity at a lower value of
  $T_{\mathrm{nuc}}$ which will go $\downarrow$. A lower temperature
  corresponds to a longer $t_{\mathrm{nuc}}$ ($\uparrow$). This would
  result in more neutrons decaying into protons and thus less
  $Y_{\mathrm{prim}}$ $\downarrow$.

  \begin{center}
  \begin{tabular}{ll}
  $\eta$ & $\downarrow$\\
  $N_{\mathrm{D}}/N_{\mathrm{n}}(T)$ & $\downarrow$\\
  $T_{\mathrm{nuc}}$ ($N_{\mathrm{D}} = N_{\mathrm{n}}$) & $\downarrow$\\
  $t_{\mathrm{nuc}}$ & $\uparrow$\\
  $X_{\mathrm{n,\,nuc}}$ & $\downarrow$\\
  $Y_{\mathrm{prim}}$ & $\downarrow$
  \end{tabular}
  \end{center}

  Each correct box fetches 0.5 marks. We give a bonus 0.5 marks if all 5
  boxes are marked correctly.
\end{description}
